% called by main.tex
%
\chapter{Diseño Experimental y Métricas de Evaluación}
\label{ch::chapter4}

Este capítulo formaliza el diseño experimental que guía todo el trabajo. El objetivo del capítulo es, definir con precisión el contrato de datos, el 
flujo de representación y el sistema de métricas con el que se evalúan las dos formulaciones
supervisadas. En otras palabras, aquí se fija qué significa que una predicción sea válida, estructuralmente fiel y útil en términos de dominio Kubernetes.

Además, el capítulo hace explícita la relación entre decisiones de representación y señales de evaluación. Esta conexión es importante porque una
salida puede parecer correcta en superficie textual y, sin embargo, fallar al reconstruir un YAML válido o incumplir requisitos de dominio. Precisamente
por eso se adopta una evaluación multinivel y no una única métrica agregada, en línea con trabajos recientes que muestran las limitaciones de evaluar
salidas estructuradas únicamente desde su parecido superficial \citep{yang2025structeval, xu2024cloudevalyaml}.

\section{Planteamiento experimental y comparabilidad}
\label{sec:ch4_design}

La pregunta experimental central es: hasta qué punto una formulación puramente secuencial puede capturar jerarquía, y qué impacto tiene sobre la
calidad estructural supervisar esa jerarquía de forma explícita en lugar de aprenderla como texto serializado, una tensión cercana a la formulación clásica de
la predicción estructurada \citep{joachims2009structured, belanger2017spen}. Para responder de forma controlada,
ambas ramas comparten modelo base, conjunto de datos, parser determinista y protocolo de evaluación.

De este modo, cualquier diferencia observada en resultados puede atribuirse a la formulación de salida y a la señal de entrenamiento estructural, no a
cambios en datos o infraestructura. De esta forma evitamos comparaciones injustas y facilitamos la interpretación de resultados.

\subsection{Formulaciones supervisadas}
\label{sec:ch4_formulations}

El estudio compara dos formulaciones:

\begin{table}[htbp]
\centering
\caption{Comparativa de formulaciones supervisadas utilizadas en el estudio.}
\label{tab:ch4_formulations}
\begin{tabular}{p{0.23\textwidth}p{0.34\textwidth}p{0.34\textwidth}}
\toprule
\textbf{Aspecto} & \textbf{Serialized SFT} & \textbf{Two-Head SFT} \\
\midrule
Superficie objetivo & Secuencia TSV con \texttt{document\_index}, \texttt{line\_index}, \texttt{level}, \texttt{line\_text}. & Secuencia TSV con contenido
y metadatos de posición; \texttt{level} fuera de la superficie textual. \\
Señal de jerarquía & Implícita en la propia cadena generada token a token. & Explícita mediante cabezal estructural dedicado para predecir \texttt{level}. \\
Parser posterior & Reconstrucción determinista a YAML, sin reparación heurística. & Reconstrucción determinista a YAML, sin reparación heurística. \\
Objetivo de la comparación & Medir hasta dónde llega la codificación superficial de estructura. & Medir ganancia al desacoplar contenido textual y jerarquía. \\
\bottomrule
\end{tabular}
\end{table}

\section{Formato de datos y representación habitual}
\label{sec:ch4_contract}

El formato de las muestras usado en este TFM parte de un principio simple: la unidad de aprendizaje no es el árbol YAML completo, sino una secuencia de bloques lineales
que conserva información suficiente para reconstruirlo de forma determinista. Cada bloque representa una línea física con metadatos estructurales.

\subsection{Unidad básica: bloque estructural}
\label{sec:ch4_block_unit}

Cada manifiesto se representa como una secuencia de tuplas de la forma:

\begin{center}
\texttt{(document\_index, line\_index, level, line\_text)}
\end{center}

donde \texttt{document\_index} identifica el documento YAML en escenarios multi-documento, \texttt{line\_index} fija el orden local de líneas en cada
documento, \texttt{line\_text} contiene el contenido sin sangría inicial y \texttt{level} codifica la indentación en unidades de dos espacios.

Una distinción clave es que \texttt{level} no equivale necesariamente a profundidad semántica del árbol parseado. Para entenderla, es útil partir de
un ejemplo concreto. Considérese el siguiente manifiesto \texttt{ConfigMap}, un tipo de recurso frecuente en el dataset de este trabajo, anotado con
el \texttt{level} que corresponde a cada línea:

\begin{lstlisting}[style=yamlblock]
apiVersion: v1       # level = 0
kind: ConfigMap      # level = 0
metadata:            # level = 0
  name: app-config   # level = 1
  labels:            # level = 1
    app: my-app      # level = 2
\end{lstlisting}

\texttt{level} cuenta unidades de sangría en el texto renderizado: cada dos espacios de indentación suman uno. La línea \texttt{app: my-app} tiene
\texttt{level = 2} porque está desplazada cuatro espacios del margen. Hasta aquí, la lectura es directa. El contraste aparece al parsear el documento
como árbol de objetos. La profundidad de un nodo se define como el número de aristas que lo separan de la raíz: la raíz tiene profundidad~0, sus
hijos directos tienen profundidad~1, y así sucesivamente. Bajo esa definición, el escalar \texttt{my-app} está a tres aristas de la raíz
(raíz $\to$ \texttt{metadata} $\to$ \texttt{labels} $\to$ \texttt{my-app}). De modo que \texttt{yaml\_depth} para este documento
vale~3, mientras que el \texttt{level} máximo observado en los bloques es~2. Las dos magnitudes suelen correlacionar, pero se calculan desde
representaciones distintas: una trabaja sobre el texto renderizado, la otra sobre la estructura del árbol parseado. En la práctica, la variable de 
\texttt{level} se corresponde con $\texttt{yaml\_depth} - 1$.   

\subsection{Granularidad del dataset durante los experimentos}
\label{sec:ch4_granularity}

Para evitar ambigüedades en la terminología, el dataset se maneja con cuatro niveles de granularidad:

\begin{itemize}
\item \textbf{Muestra} (\texttt{sample}): unidad conceptual asociada a una petición en lenguaje natural y su salida esperada.
\item \textbf{Documento} (\texttt{document}): subunidad YAML dentro de una muestra, indexada por \texttt{document\_index}.
\item \textbf{Línea/Bloque} (\texttt{line} o \texttt{block}): fila estructural con \texttt{line\_index}, \texttt{level} y \texttt{line\_text}.
\item \textbf{Nivel de línea} (\texttt{line level}): profundidad de indentación de una línea concreta, usada tanto en entrenamiento como en evaluación estructural.
\end{itemize}

Esta jerarquía de granularidad permite construir métricas locales (por línea), métricas de consistencia (por documento) y métricas de éxito (por muestra).
En la práctica, las métricas se van a construir a nivel de muestra o a nivel de línea, dependiendo de la señal que se quiera capturar. Por ejemplo, la tasa de parseo
es una métrica de muestra (¿parsea el documento completo?), mientras que la exactitud de nivel es una métrica de línea (¿predice el nivel correcto para esta línea?).

\subsection{Pipeline de transformación: YAML a CSV/TSV}
\label{sec:ch4_pipeline}

El flujo de datos habitual durante los experimentos sigue cuatro etapas:

\begin{enumerate}
\item Normalización de YAML de referencia para fijar una forma canónica estable.
\item Conversión de YAML canónico a secuencia de bloques estructurados.
\item Serialización de bloques a formato tabular de entrenamiento (principalmente TSV, con exportación CSV/TSV para análisis externo).
\item Reconstrucción determinista de YAML desde bloques para validación \textit{round-trip}.
\end{enumerate}

La elección de una representación tabular se justifica por su versatilidad. Permite entrenar con una interfaz textual homogénea, inspeccionar errores por columna y trazar
fallos. Además, al mantener \texttt{document\_index} y \texttt{line\_index}, la evaluación conserva orden y segmentación.

El parser cumple aquí una función de control estructural: delimita qué salidas pueden reconstruirse de forma segura. Esta idea se relaciona con trabajos
de generación guiada y parsers incrementales, donde la validez formal de la salida se incorpora como frontera explícita del proceso de generación o
validación \citep{scholak2021picard, geng2023grammar, willard2023efficient}.

\begin{table}[htbp]
\centering
\caption{Serializaciones operativas usadas en el proyecto.}
\label{tab:ch4_serializations}
\begin{tabular}{p{0.28\textwidth}p{0.26\textwidth}p{0.38\textwidth}}
\toprule
\textbf{Formato} & \textbf{Columnas visibles} & \textbf{Uso principal} \\
\midrule
\texttt{blocks\_tsv\_v1} & \texttt{document\_index}, \texttt{line\_index}, \texttt{level}, \texttt{line\_text} & Supervisión secuencial estructurada en Serialized SFT. \\
\texttt{blocks\_tsv\_compact\_v1} & \texttt{document\_index}, \texttt{level}, \texttt{line\_text} & Variante compacta para análisis y baselines controlados. \\
\texttt{content\_blocks\_v1} & \texttt{document\_index}, \texttt{line\_index}, \texttt{line\_text} & Superficie textual para Two-Head SFT, con \texttt{level}
como etiqueta separada. \\
\bottomrule
\end{tabular}
\end{table}

\begin{figure}[htbp]
\centering
\fbox{%
\parbox{0.9\textwidth}{
\textbf{Esquema del pipeline (visión conceptual)}\\
YAML canónico $\rightarrow$ Bloques estructurados $\rightarrow$ Serialización tabular (TSV/CSV) $\rightarrow$ Predicción del modelo $\rightarrow$ Reconstrucción YAML $\rightarrow$ Evaluación multinivel.
}}
\caption{Pipeline metodológico desde la representación canónica hasta la evaluación.}
\label{fig:ch4_pipeline_scheme}
\end{figure}

\section{Protocolo de entrenamiento y evaluación}
\label{sec:ch4_protocol}

La comparativa experimental se apoya en tres principios: mismo split de datos entre formulaciones, misma base de modelo y mismo mecanismo de
reconstrucción y evaluación. El modelo base es Qwen2.5-7B-Instruct \citep{qwen2024qwen25} en cuantización de 4 bits. La cuantización consiste en representar los pesos
del modelo con una precisión reducida (cuatro bits por parámetro frente a los 16 o 32 habituales), lo que reduce de forma
importante el consumo de memoria de la GPU sin comprometer en exceso la calidad de las representaciones internas \citep{dettmers2023qlora}. Para adaptar el modelo al 
hardware disponible, el ajuste fino se
realiza mediante LoRA (\textit{Low-Rank Adaptation}), una técnica que introduce matrices de adaptación de bajo rango (rango 8 en este
trabajo) en las proyecciones de atención del modelo. Esto permite modificar su comportamiento sin tocar los pesos originales ni requerir los
recursos de un entrenamiento completo \citep{hu2021lora}.

En este trabajo la principal diferencia entre ramas es la forma en la que se representa y
supervisa la jerarquía. Debe recordarse que se parte de la hipótesis de que los errores de generación estructurada no se explican solo
por semántica de contenido, sino también por cómo la arquitectura aprende la señal de profundidad.

\section{Métricas de evaluación: familias y utilidad}
\label{sec:ch4_metrics}

La evaluación se organiza en familias complementarias. Ninguna métrica individual basta para caracterizar calidad estructural en YAML/Kubernetes \citep{chen2021evaluating, yang2025structeval, xu2024cloudevalyaml}. Por eso, el
análisis debe combinar la validez sintáctica, fidelidad de estructura, adecuación textual y consistencia de dominio si quiere ser útil como medida cuantitativa.

\begin{table}[htbp]
\centering
\caption{Familias de métricas y pregunta metodológica que responden.}
\label{tab:ch4_metric_families}
\begin{tabular}{p{0.27\textwidth}p{0.27\textwidth}p{0.37\textwidth}}
\toprule
\textbf{Familia} & \textbf{Ejemplos} & \textbf{Utilidad metodológica} \\
\midrule
Sintácticas de parseo & parseo de salida estructurada, parseo YAML, contrato de bloques & Determinar si la salida es procesable sin intervención manual. \\
Estructurales de jerarquía & exactitud y MAE de nivel, recall profundo 5--8 & Medir fidelidad de la geometría de indentación y orden interno. \\
Fidelidad textual & exactitud de contenido de línea, precisión/recobrado/F1 textual, BLEU/ROUGE auxiliares & Cuantificar cuánto contenido de referencia se conserva. \\
Semántica Kubernetes & F1 de claves semánticas, coincidencia de \texttt{kind}, \texttt{apiVersion}, selectores/volúmenes & Verificar corrección funcional básica de
recursos generados. \\
Requisitos de prompt y campos requeridos & precisión/recobrado de requisitos, campos obligatorios completos por muestra & Evaluar alineación entre petición natural y
manifiesto resultante. \\
Validez de dominio multinivel & score KDV, nivel de validez 0--5, tasa de paso de compuertas de dominio & Integrar señales estructurales y semánticas en un criterio operativo único. \\
\bottomrule
\end{tabular}
\end{table}

\subsection{Métricas sintácticas y estructurales}
\label{sec:ch4_metrics_structural}

Las métricas de parseo son la primera barrera: sin parseo correcto, no tiene sentido discutir calidad semántica.

Este bloque de métricas responde a la pregunta: \textit{\textquestiondown la salida es internamente coherente como objeto estructurado?}

Las principales métricas de este bloque empiezan antes del propio YAML reconstruido. En primer lugar se registra si la salida textual del modelo puede convertirse en una secuencia
de bloques válidos. Si \(S\) es el conjunto de filas generadas y \(E \subseteq S\) el subconjunto que supera el parseo de salida estructurada, la tasa
\texttt{structured\_output\_parse\_success\_rate} se define como:

\[
\texttt{structured\_output\_parse\_success\_rate}
= \frac{|E|}{|S|}.
\]

Sobre ese subconjunto evaluable se calcula después la tasa de parseo YAML correcto. Si \(\mathrm{yaml}(y_i)\) indica que la reconstrucción determinista de la predicción \(y_i\)
puede cargarse como YAML sin error, entonces:

\[
\texttt{yaml\_parse\_success\_rate}
= \frac{1}{|E|}\sum_{i \in E}\mathbf{1}[\mathrm{yaml}(y_i)].
\]

Esta separación es importante porque una salida puede respetar la superficie tabular esperada y, aun así, producir una trayectoria de niveles incompatible con un YAML parseable.

\subsection{Métricas estructurales de jerarquía}

Una vez superada esa primera frontera, se calculan las métricas de fidelidad estructural por línea. De estas últimas, las dos más informativas son
la tasa de coincidencia exacta de nivel (\texttt{level\_exact\_match\_rate}) y el error absoluto medio de nivel (\texttt{level\_mae}). Sean $R = (r_1, \dots, r_n)$
los niveles de referencia y $P = (p_1, \dots, p_m)$ los niveles predichos, y sea $k = \min(n, m)$ el número de líneas comparables:

$$\texttt{level\_exact\_match\_rate} = \frac{1}{n}\sum_{i=1}^{k} \mathbf{1}[r_i = p_i]$$

$$\texttt{level\_mae} = \frac{1}{k}\sum_{i=1}^{k} |r_i - p_i|$$

Ambas se calculan solo sobre el prefijo común de longitud~$k$: si la predicción genera menos líneas que la referencia, las líneas no cubiertas no
cuentan como coincidencias pero tampoco penalizan directamente el MAE. Esa elección de implementación implica que errores de recuento de líneas se
capturan mejor con \texttt{line\_count\_match} que con \texttt{level\_mae}.

Para los experimentos con cabezal explícito de nivel se introduce además una métrica centrada en la región profunda de la jerarquía, porque los niveles superficiales dominan el
promedio global. Sea \(D = \{i \leq k \mid r_i \in \{5,6,7,8\}\}\) el conjunto de posiciones comparables cuyo nivel real pertenece a la zona profunda. El recall exacto sobre
esa región es:

\[
\texttt{deep\_level\_exact\_recall\_5\_8}
=
\frac{\sum_{i=1}^{k}\mathbf{1}[r_i \in \{5,6,7,8\}]\,\mathbf{1}[p_i = r_i]}
{\sum_{i=1}^{k}\mathbf{1}[r_i \in \{5,6,7,8\}]}.
\]

Si no existen líneas de referencia en niveles 5--8, el soporte profundo es cero y la implementación registra el recall como \(0\). Su función no es sustituir a la exactitud
global, sino impedir que un modelo parezca estructuralmente sólido solo porque acierta los niveles más frecuentes.

\subsection{Métricas de fidelidad textual y semántica}
\label{sec:ch4_metrics_semantic}

Una vez garantizada la estructura mínima, se analiza contenido. En este punto conviene precisar el papel de las métricas de solapamiento textual como
BLEU o ROUGE. Estas métricas, ampliamente utilizadas en tareas de generación de lenguaje natural como traducción automática o resumen, miden en qué
medida la salida comparte n-gramas con una referencia. En un dominio de texto libre esa señal es razonablemente informativa. En el caso de manifiestos
Kubernetes, sin embargo, resulta insuficiente como criterio principal: una salida puede conservar la mayoría de los tokens de referencia y aun así
producir un YAML que no parsea, porque una sola línea con la indentación incorrecta rompe la estructura del documento. A la inversa, una salida
estructuralmente válida pero que emplea nombres de campo o valores distintos puede tener un ROUGE alto sin ser funcionalmente equivalente. Por eso
BLEU y ROUGE se usan aquí en un rol auxiliar, como señal secundaria de solapamiento de contenido, pero no como criterio de corrección estructural
ni de validez semántica en el dominio, una distinción coherente con la evaluación funcional de artefactos técnicos y con benchmarks específicos de
YAML cloud-native \citep{chen2021evaluating, xu2024cloudevalyaml}.

Dentro de las métricas de fidelidad textual sí intervienen directamente en el análisis la precisión, el recobrado y el F1 de contenido de línea.
Sean $\hat{T}$ el multiconjunto de líneas de texto predichas y $T$ el de referencia; los verdaderos positivos son el solapamiento multiconjunto
entre ambos, $TP = |\hat{T} \cap T|$:

$$\text{Precision} = \frac{TP}{|\hat{T}|}, \qquad \text{Recall} = \frac{TP}{|T|}, \qquad F_1 = \frac{2 \cdot \text{Precision} \cdot \text{Recall}}{\text{Precision} + \text{Recall}}$$

donde $|\cdot|$ denota cardinalidad con multiplicidades. Esta familia es sensible a qué contenido se genera, pero ciega al orden de las líneas
dentro del manifiesto, que es precisamente lo que capturan las métricas estructurales del bloque anterior.

Las métricas semánticas de Kubernetes, por su parte, evalúan si la salida conserva campos críticos y relaciones internas relevantes: coherencia entre
selectores y etiquetas, correspondencia entre \texttt{volumeMounts} y \texttt{volumes}, presencia y corrección de \texttt{kind} y \texttt{apiVersion}.
Estas señales apuntan a si el manifiesto generado es funcionalmente plausible, algo que ninguna métrica puramente textual puede capturar.

La métrica \texttt{average\_semantic\_key\_f1} resume una parte de esa señal mediante presencia de claves. Sea \(\mathcal{K}\) el vocabulario fijo de claves semánticas relevantes
para el dominio (por ejemplo \texttt{metadata}, \texttt{spec}, \texttt{containers}, \texttt{image}, \texttt{ports}, \texttt{volumes}, \texttt{volumeMounts},
\texttt{selector} o \texttt{template}). Para una muestra \(i\), sean \(K_i^\star \subseteq \mathcal{K}\) las claves presentes en el YAML de referencia y
\(\widehat{K}_i \subseteq \mathcal{K}\) las claves presentes en el YAML predicho. Entonces:

\[
\mathrm{SemPrec}_i = \frac{|K_i^\star \cap \widehat{K}_i|}{|\widehat{K}_i|}, \qquad
\mathrm{SemRec}_i = \frac{|K_i^\star \cap \widehat{K}_i|}{|K_i^\star|},
\]

\[
\mathrm{SemF1}_i =
\frac{2 \cdot \mathrm{SemPrec}_i \cdot \mathrm{SemRec}_i}
{\mathrm{SemPrec}_i + \mathrm{SemRec}_i},
\qquad
\texttt{average\_semantic\_key\_f1}
= \frac{1}{|E|}\sum_{i \in E}\mathrm{SemF1}_i.
\]

Las divisiones con denominador cero se resuelven con la convención de implementación \(\mathrm{safe\_divide}(a,0)=0\). En la práctica, esto evita que una muestra sin claves
semánticas predichas produzca valores indefinidos.

Aquí la pregunta cambia: \textit{¿además de estar bien formado, el manifiesto representa lo que debe representar?} La combinación de señales
textuales y semánticas evita dos falsos positivos habituales: salidas gramaticalmente válidas pero funcionalmente incorrectas, y salidas
semánticamente plausibles pero estructuralmente frágiles.

\subsection{Métricas de requisitos de prompt y campos requeridos}
\label{sec:ch4_metrics_prompt}

El sistema también evalúa la alineación con requisitos extraídos del prompt, una dimensión especialmente relevante en tareas que traducen intención natural a manifiestos
técnicos \citep{ueno2024migrating, angi2025kgen}. La idea es que una predicción puede estar bien formada semánticamente
y aun así ignorar restricciones concretas del usuario: el tipo de recurso pedido, el nombre, la imagen del contenedor o el número de réplicas.
Para capturar esa brecha, el sistema extrae automáticamente un conjunto de requisitos~$Q$ desde el texto del prompt y un conjunto de
campos cubiertos~$C$ desde el YAML predicho. La precisión y el recobrado de requisitos se definen entonces como:

$$\text{Req-Precision} = \frac{|Q \cap C|}{|C|}, \qquad \text{Req-Recall} = \frac{|Q \cap C|}{|Q|}$$

Cuando $|Q| = 0$, es decir, el prompt no contiene requisitos extraibles; la métrica se marca como no aplicable y no entra en los promedios
del dataset. Cuando $|Q| = |C|$ y el conjunto coincide exactamente, la muestra cuenta como \texttt{prompt\_requirement\_exact\_match}.

Por otro lado, el evaluador mide si los recursos generados contienen sus campos obligatorios mínimos. Sea \(R_i\) el conjunto de recursos Kubernetes predichos en la muestra \(i\),
y sea \(F(\mathrm{kind}(r))\) el conjunto de campos requeridos para el tipo de recurso \(r\). Definimos primero la completitud de un recurso como:

\[
\mathrm{complete}(r) =
\prod_{f \in F(\mathrm{kind}(r))}\mathbf{1}[f \in r].
\]

La muestra completa todos sus campos obligatorios si todos sus recursos aplicables son completos:

\[
\mathrm{RequiredCompleteSample}(y_i) =
\prod_{r \in R_i}\mathrm{complete}(r).
\]

Agregada sobre el conjunto de muestras evaluadas con recursos aplicables, esta señal produce:

\[
\texttt{required\_field\_complete\_sample\_rate}
=
\frac{1}{|E_\mathrm{req}|}
\sum_{i \in E_\mathrm{req}}
\mathrm{RequiredCompleteSample}(y_i).
\]

Esta métrica es deliberadamente más estricta que una tasa de presencia de campos calculada recurso a recurso: una sola omisión en una muestra multi-recurso hace que la muestra no
cuente como completa. Por eso se puede usar como indicador de integridad mínima del manifiesto completo.

\subsection{Validez de dominio}
\label{sec:ch4_metrics_domain}

Las métricas anteriores responden a preguntas locales: si una línea tiene el nivel correcto, si el YAML parsea, si cierto campo está presente.
Ninguna de ellas, por sí sola, dice si un manifiesto es operativamente aceptable. Para eso el sistema define una escala de validez de dominio
estructurada en seis compuertas secuenciales, numeradas del 0 al 5, que se evalúan en orden y se detienen en la primera que falla.

La idea es imitar la cadena de comprobaciones que un operador real aplicaría antes de aplicar un manifiesto a un clúster, teniendo en cuenta que los defectos de configuración no se reducen
a errores de sintaxis y pueden requerir validaciones progresivas \citep{zhang2025kubernetesdefects, kakarla2024diffy}: primero que el YAML sea parseable,
luego que la representación interna sea coherente, después que el recurso sea reconocible como Kubernetes, y así progresivamente hasta llegar a
requisitos de calidad estática. El nivel que alcanza la predicción antes de fallar (o 5 si supera todas las compuertas) es el
\texttt{kubernetes\_domain\_validity\_level}.

\begin{table}[htbp]
\centering
\caption{Escala de validez de dominio Kubernetes: compuertas secuenciales.}
\label{tab:ch4_domain_levels}
\begin{tabular}{p{0.06\textwidth}p{0.25\textwidth}p{0.58\textwidth}}
\toprule
\textbf{Nivel} & \textbf{Compuerta} & \textbf{Condición} \\
\midrule
0 & Parseo YAML & El texto generado es YAML válido y puede cargarse sin error de sintaxis. \\
1 & Contrato de bloques & Los bloques del parser satisfacen el contrato estructural completo: índices monótonos, sin fugas de indentación y ratio de bloques válidos igual a 1. \\
2 & Identidad Kubernetes & Cada documento tiene \texttt{apiVersion}, \texttt{kind} reconocido y \texttt{metadata.name} no vacíos, y los campos obligatorios del tipo de recurso están presentes. \\
3 & Invariantes intra-recurso & Coherencia dentro de cada recurso: selectores con sus etiquetas de template, imágenes de contenedor presentes, puertos en rango válido, \texttt{volumeMounts} con su \texttt{volume} correspondiente. \\
4 & Invariantes inter-recurso & Referencias entre recursos: selector de \texttt{Service} apunta a etiquetas de un workload local, \texttt{RoleBinding} apunta a un \texttt{Role} existente, \texttt{ServiceAccount} referenciada existe en el manifiesto. \\
5 & Calidad estática & Ausencia de olores de seguridad: sin imagen \texttt{latest}, \texttt{runAsNonRoot: true}, \texttt{readOnlyRootFilesystem: true}, recursos de CPU y memoria declarados, sin namespaces de host habilitados. \\
\bottomrule
\end{tabular}
\end{table}

La escala es acumulativa: superar el nivel~$\ell$ implica haber superado todos los anteriores. Si \(c_r(y_i)\) representa que la muestra \(i\) supera la compuerta \(r\),
la tasa de paso del nivel~\(\ell\) sobre un conjunto de evaluación \(S\) puede escribirse como:

\[
\mathrm{pass}_{\ell} =
\frac{1}{|S|}\sum_{i \in S}\prod_{r=0}^{\ell}\mathbf{1}[c_r(y_i)].
\]

Además del nivel alcanzado y de la tasa de paso de cada compuerta, el evaluador conserva una señal graduada denominada \texttt{kubernetes\_domain\_validity\_score}. Para cada
muestra se construye un vector binario de seis componentes \(Q_i = (q_{i,0}, \ldots, q_{i,5})\), donde \(q_{i,r}=1\) si la comprobación individual asociada al nivel \(r\) se satisface
y \(q_{i,r}=0\) en caso contrario. El score KDV de la muestra y su media agregada se definen como:

\[
\mathrm{KDV}(y_i) =
\frac{1}{6}\sum_{r=0}^{5} q_{i,r},
\qquad
\overline{\mathrm{KDV}}
=
\frac{1}{|S|}\sum_{i \in S}\mathrm{KDV}(y_i).
\]

La media \(\overline{\mathrm{KDV}}\) se registra en los experimentos como \texttt{average\_\allowbreak kubernetes\_\allowbreak domain\_\allowbreak validity\_\allowbreak score}.
Esta magnitud es más suave que \texttt{kubernetes\_\allowbreak domain\_\allowbreak gate\_\allowbreak pass\_\allowbreak rate}: una muestra que no alcanza el nivel~5 puede seguir
recibiendo crédito parcial por las comprobaciones que sí satisface, aunque el nivel acumulativo se detenga en la primera compuerta fallida. Precisamente por eso resulta útil en comparaciones
intermedias y en la construcción de preferencias automáticas, donde muchas candidatas pueden fallar la compuerta final pero diferir claramente en calidad de dominio.

La definición acumulativa de \(\mathrm{pass}_{\ell}\) convierte la escala en un criterio legible de comparación
entre formulaciones, más informativo que una media de métricas heterogéneas. Si un modelo pasa sistemáticamente el nivel~2 pero no el~3, el diagnóstico
es claro: genera recursos identificables pero con incoherencias internas. Si se queda en el nivel~0, el problema es más básico: no produce YAML parseable.

La compuerta final, el nivel~5, recoge requisitos de calidad que no son estrictamente obligatorios para que Kubernetes acepte el manifiesto, pero que
sí son buenas prácticas habituales en entornos de producción. Precisamente por eso es la más difícil de superar y la que actúa como criterio de
excelencia operativa, especialmente porque las configuraciones Kubernetes pueden ser sintácticamente válidas y aun así introducir riesgos de seguridad o
mantenimiento \citep{rahman2023securityk8s, zhang2024generativeSmells}. La tasa de muestras que alcanzan el nivel~5 se registra como
\texttt{kubernetes\_\allowbreak domain\_\allowbreak gate\_\allowbreak pass\_\allowbreak rate}, y será
uno de los indicadores principales en el capítulo de resultados.

Finalmente, todas las métricas se agregan a nivel de dataset para obtener tasas medias, varianzas y criterios de comparación entre formulaciones.
Sea $S$ el conjunto de muestras evaluadas; para cualquier métrica binaria $x_i \in \{0,1\}$ la tasa agregada es simplemente
$\bar{x} = \frac{1}{|S|}\sum_{i \in S} x_i$, y para métricas continuas como F1 o MAE se calcula la media aritmética excluyendo muestras
no aplicables. Esta agregación permite pasar de diagnósticos locales a conclusiones experimentales robustas, que se desarrollarán en el capítulo de resultados.
